% Flavor catalog per-process entry.
% process_id: CR007
% family: collider_rs
% owner_agent_id: PKA-CR007
% writer_agent_id: WA
% checker_agent_id: null
% opus_signoff_id: null
% Canonical metadata sidecar: CR007.yaml
% Status is controlled by the YAML sidecar status_history, not this comment.

\section{CR007: Spin-2 KK-Graviton Resonance in Bulk RS}

\subsection*{Process}
\textbf{Standard notation:}
\[
  pp \to G^{(1)}_{\rm KK} \to ZZ,\ WW,\ \gamma\gamma,\ \ell\ell,\ t\bar t .
\]
This entry covers direct LHC searches for a heavy spin-2 Kaluza--Klein
graviton in Randall--Sundrum models, with emphasis on the bulk-RS benchmark in
which Standard Model fields propagate in the warped extra dimension.  The
experimentally quoted final states are narrow or moderately broad resonance
searches in diboson, diphoton, dilepton, and top-pair spectra; the RS
interpretation maps the null result onto a lower limit on the first graviton
mode mass for a specified curvature coupling \(k/\overline{M}_{\rm Pl}\).

\subsection*{Current best limit(s)}
The PDG 2025 Extra Dimensions listing is the canonical one-place summary in
\texttt{flavor\_catalog/references/CR007/pdg2025\_extra\_dimensions\_rsg\_extract.txt}.
For the generic RS-graviton benchmark at \(k/\overline{M}_{\rm Pl}=0.1\), PDG
lists the strongest current mass lower limit as
\[
  m(G^{(1)}_{\rm KK}) > 4.8~{\rm TeV}\quad(95\%~{\rm CL}),
\]
from the CMS high-mass diphoton search, with the CMS high-mass dilepton result
giving the adjacent \(m(G^{(1)}_{\rm KK})>4.78~{\rm TeV}\) entry.  These
headline limits are powerful generic spin-2 resonance bounds, but they are not
the most faithful bulk-RS interpretation when light-fermion and photon
couplings are suppressed.

For the bulk-RS diboson benchmark, the current PDG-listed full Run-2 CMS
all-jets search gives
\[
  m(G_{\rm bulk}) > 1.4~{\rm TeV}\quad(95\%~{\rm CL})
\]
for \(k/\overline{M}_{\rm Pl}=0.5\) in \(WW/ZZ\) final states.  The earlier
ATLAS combined bosonic-plus-leptonic search at \(36.1~{\rm fb}^{-1}\) excluded
a bulk-RS KK graviton below \(2.3~{\rm TeV}\) for its \(k/\overline{M}_{\rm Pl}=1\)
benchmark.  These mass numbers sit in the sidecar
\texttt{pdg\_or\_equivalent.values} block; luminosities, production-mode
assumptions, and benchmark-coupling choices are recorded as supporting or
auxiliary inputs rather than promoted measured observables.

\subsection*{Relevance to RS}
Bulk RS graviton searches constrain a direct collider resonance mass, not the
\(\Delta F=2\) Wilson coefficients that dominate the existing quark-flavor
pipeline.  They are nevertheless useful because custodial bulk-RS spectra often
contain correlated KK gauge bosons, KK gravitons, vectorlike fermions, and
radion-like states at nearby TeV scales.  The spin-2 channel is therefore an
RS-distinguishing cross-check of any parameter point that places the first
graviton in the LHC range.

The current low-energy quark scan gives
\(M_{\rm KK}^{\min}(p50,\ g_\ast=3)=47.26~{\rm TeV}\) in
\texttt{docs/quark\_scan\_methodology\_note.tex}.  Compared with that
anarchic-flavor bound, the LHC graviton limits above are far weaker as mass
setters.  Their value is instead diagnostic: they can veto or reinterpret
non-anarchic, flavor-protected, or custodial RS scenarios whose flavor
constraints have been deliberately weakened, and they provide an external check
on whether the scan has wandered into spectra that would already have produced
a high-mass diboson, diphoton, dilepton, or \(t\bar t\) excess.

\subsection*{Post-2010 developments}
The original bulk-RS graviton collider expectation in
\texttt{AgasheDavoudiaslPerezSoni2007\_WarpedGravitons} emphasized suppressed
light-fermion and photon channels but enhanced sensitivity in longitudinal
\(W/Z\) modes.  LHC searches after 2010 tested that pattern in stages.  Early
Run-1 searches established the first ATLAS/CMS mass and cross-section limits in
\(ZZ\), \(WW\), and combined \(\gamma\gamma/\ell\ell\) channels, replacing the
Tevatron reach.  The ATLAS \(36.1~{\rm fb}^{-1}\) combination
(\texttt{arXiv:1808.02380}) mattered because it combined bosonic and leptonic
Run-2 final states into a single bulk-RS spin-2 interpretation.

The full Run-2 ATLAS diboson searches (\texttt{arXiv:1906.08589} and
\texttt{arXiv:2004.14636}) moved the diboson program into all-hadronic and
semileptonic boosted topologies with \(139~{\rm fb}^{-1}\).  CMS then provided
the strongest generic RS spin-2 dilepton benchmark
(\texttt{arXiv:2103.02708}), the current bulk-graviton all-jets diboson
benchmark used here (\texttt{arXiv:2210.00043}), and the current PDG-leading
diphoton benchmark (\texttt{arXiv:2405.09320}).  The progression is mainly
in luminosity, boosted-boson tagging, and channel coverage; it is not a move to
a model-independent RS likelihood.

\subsection*{Validity / model dependence}
The exclusion limits are conditional.  They assume a production mechanism,
resonance width, \(k/\overline{M}_{\rm Pl}\), acceptance, and decay-branching
pattern.  A headline \(G\to\gamma\gamma\) or \(G\to\ell\ell\) limit is not
automatically applicable to bulk RS models where light-fermion and photon
couplings are suppressed.  Conversely, a \(G_{\rm bulk}\to WW/ZZ\) limit does
not directly constrain the KK gluon mass or a custodial fermion mass without a
specified spectrum and branching matrix.

For the anarchic-flavor catalog, the honest use is therefore as a direct-search
side condition: compare the predicted \(G^{(1)}_{\rm KK}\) mass, width, and
branching fractions against the relevant ATLAS/CMS limit curve.  A single
mass-threshold cut is only defensible for the exact experimental benchmark
quoted in the sidecar.

\subsection*{Code coverage}
\textbf{NO.}  Greps over \texttt{quarkConstraints/}, \texttt{qcd/},
\texttt{flavorConstraints/}, \texttt{neutrinos/}, \texttt{yukawa/},
\texttt{warpConfig/}, \texttt{solvers/}, \texttt{scanParams/}, and
\texttt{tests/} found no direct ATLAS/CMS, diboson, diphoton, dilepton,
graviton-resonance, or collider-reinterpretation implementation.  The only
``CMS'' hit in the collider grep is an unrelated RunDec calibration comment in
\texttt{tests/test\_alpha\_s.py:89}.  Adjacent code does carry the scan's
flavor KK scale: \texttt{quarkConstraints/modern/scan.py:1229--1255}
enumerates \(M_{\rm KK}\) for flavor scan points, and
\texttt{quarkConstraints/modern/evaluation.py:643--666} evaluates the
\(\Delta F=2\) flavor matching at that scale.  No separate collider-direct
filter is applied.

\subsection*{Implementation difficulty}
\textbf{HIGH.}  A faithful implementation needs more than a scalar mass cut:
it needs the model point's spin-2 graviton mass, width, production cross
section, branching ratios into \(WW\), \(ZZ\), \(\gamma\gamma\), \(\ell\ell\),
and \(t\bar t\), plus acceptance or a recast of the experimental analysis.
That is naturally an external reinterpretation workflow using tools such as
MadGraph/Pythia/Delphes with CheckMATE, MadAnalysis5, Rivet, or a HEPData-based
likelihood approximation.  The repo currently has neither collider event
simulation nor an RS-resonance likelihood interface.

\subsection*{Key references}
Process-local source keys before bibliography consolidation are
\texttt{PDG2025\_ExtraDimensions\_RSG}, \texttt{CMS2024\_Diphoton},
\texttt{CMS2023\_AllJetsBosonPairs}, \texttt{CMS2021\_Dilepton},
\texttt{ATLAS2020\_SemileptonicDiboson},
\texttt{ATLAS2019\_HadronicDiboson}, \texttt{ATLAS2018\_Combination},
\texttt{AgasheDavoudiaslPerezSoni2007\_WarpedGravitons},
\texttt{AgasheEtAl2007\_WarpedGaugeBosons}, and
\texttt{RandallSundrum1999\_Hierarchy}.
